\chapter{Ηλεκτρικά Χαρακτηριστικά Εδάφους}
\label{ch:electrical}

\section{Εισαγωγή}

Στο παρόν κεφάλαιο παρουσιάζεται η πειραματική διαδικασία που
ακολουθήθηκε για τον ηλεκτρικό χαρακτηρισμό των δειγμάτων εδάφους μέσω
Φασματοσκοπίας Εμπέδησης (Electrochemical Impedance Spectroscopy, EIS).
Αρχικά περιγράφονται οι βασικές αρχές της τεχνικής και η χρησιμότητά της
για τον χαρακτηρισμό του συστήματος. Στη συνέχεια, παρουσιάζεται η
μεθοδολογία που εφαρμόστηκε καθώς και τα συμπεράσματα για το ηλεκτρικό
ισοδύναμο του δείγματος εδάφους που εξετάστηκε.

\section{Φασματοσκοπία Εμπέδησης}

Η Φασματοσκοπία Εμπέδησης είναι μία τεχνική που εφαρμόζει ένα
ημιτονοειδές εναλλασσόμενο ηλεκτρικό σήμα μικρού πλάτους σε ένα υπό
εξέταση σύστημα και μετρά την απόκρισή του σε ένα φάσμα συχνοτήτων. Όταν
το δείγμα διεγείρεται, για κάθε διαφορετικό εύρος συχνοτήτων κυριαρχούν
και οι αντίστοιχοι μηχανισμοί. Η απόκριση του συστήματος εκφράζεται μέσω
της σύνθετης αντίστασης και μπορεί να περιγράψει τους μηχανισμούς αυτούς
μοντελοποιώντας το σύστημα με ένα αντίστοιχο ηλεκτρικό κυκλωματικό
ισοδύναμο. \cite{ref22}

\section{Εξοπλισμός και Μέθοδος}

Για τη μελέτη της απόκρισης των δειγμάτων συναρτήσει της συχνότητας
χρησιμοποιήθηκε ο Αναλυτής Σύνθετης Εμπέδησης (Impedance Analyzer) της
Keysight Technologies E4990A. Το εύρος συχνοτήτων που αξιολογήθηκαν
είναι από 20Hz έως 1.5MHz, συχνότητες που θεωρήθηκαν ασφαλείς μετά τη
δοκιμή γνωστών στοιχείων. Με στόχο τη μέγιστη ακρίβεια, η βαθμονόμιση
του οργάνου και οι μετρήσεις έγιναν τμηματικά στο φάσμα της συχνότητας,
προκειμένου να ελαχιστοποιηθούν τα συχνοεξαρτώμενα σφάλματα του αναλυτή
και των καλωδίω, και με στόχο να αυξηθεί ο αριθμός των συνολικών
δειγμάτων. Η τάση που εφαρμόζει ο αναλυτής είναι 500mVrms, τάση που
οφείλει να είναι χαμηλή ώστε να μην υπερφορτώνεται το δείγμα. Ο τύπος
του δείγματος που χρησιμοποιήθηκε εντάσσεται στην κατηγορία Sandy Loam.
\cite{ref31}

\section{Απόκριση Συστήματος}

Η απόκριση του συστήματος για (f)=20Hz-1.5MHz και gwc 10\% φαίνεται στο
\figref{fig:impedance-magnitude}. Είναι χαρακτηριστική η μεγάλη τιμή του μέτρου της εμπέδησης (
\textbar Ζ\textbar{} ) στις πολύ χαμηλές συχνότητες, γεγονός που μπορεί
να δικαιολογηθεί από την πόλωση των ηλεκτροδίων. Από τα 10kHz και μετά,
παρατηρείται σταδιακή μείωση του \textbar Z\textbar. Ωστόσο, στις
υψηλότερες συχνότητες το \textbar Z\textbar{} τείνει να σταθεροποιηθεί,
γεγονός που υποδηλώνει ότι στην περιοχή αυτή το \textbar Z\textbar{}
επηρεάζεται λιγότερο από την εφαρμοζόμενη συχνότητα.

Η πόλωση των ηλεκτροδίων είναι ένα φαινόμενο που εμφανίζεται μεταξύ
ηλεκτροδίου-μέσου όταν εφαρμόζεται ηλεκτρομαγνητικό πεδίο και δεν
υπάρχει πραγματική μεταφορά φορτίου. Τα φορτία συσσωρεύονται αλλά δεν
περνούν στο μέσο, σχηματίζοντας τη λεγόμενη Ηλεκτρική Διπλοστιβάδα
(Electrical Double Layer) η οποία προκαλεί τη μεγάλη τιμή του
\textbar Z\textbar{} στις χαμηλές συχνότητες (\figref{fig:electrode-polarization}).

\begin{figure}[H]
  \centering
  \includegraphics[width=0.64\textwidth]{impedance-magnitude-vs-frequency.png}
  \caption{Συχνοτική απόκριση δείγματος με gwc 10\% από 20Hz έως 1,5MHz.}
  \label{fig:impedance-magnitude}
\end{figure}

\begin{figure}[H]
  \centering
  \includegraphics[width=0.66\textwidth]{electrode-polarization.png}
  \caption{Το φαινόμενο της Ηλεκτρικής Διπλοστιβάδας μεταξύ ηλεκτροδίου και μέσου. \cite{ref30}}
  \label{fig:electrode-polarization}
\end{figure}

Στο \figref{fig:phase-frequency} φαίνεται το διάγραμμα της φάσης συναρτήσει της συχνότητας.
Στις χαμηλές συχνότητες εμφανίζεται ξανά η πόλωση των ηλεκτροδίων, ενώ
παραλείποντας αυτό το εύρος, η φάση βρίσκεται πολύ κοντά στο μηδέν και
με σταδιακή αύξηση της συχνότητας γίνεται αρνητικότερη.

\begin{figure}[H]
  \centering
  \includegraphics[width=0.64\textwidth]{phase-vs-frequency.png}
  \caption{Διάγραμμα φάσης-Συχνότητας για δείγμα με gwc 10\% από 20Hz έως 1,5MHz.}
  \label{fig:phase-frequency}
\end{figure}

Ένα χρήσιμο εργαλείο που χρησιμοποιείται συχνά στην Φασματοσκοπία
Εμπέδησης είναι το διάγραμμα Nyquist που φαίνεται στο \figref{fig:nyquist-full}. Η
συχνότητα αυξάνεται από τα δεξιά προς τα αριστερά, στον Χ-άξονα
απεικονίζεται η πραγματική συνιστώσα Z' που εκφράζει την ωμική
συμπεριφορά του συστήματος και στον Υ-άξονα η φανταστική συνιστώσα
-Ζ'\,' που εκφράζει είτε τη χωρητική είτε την επαγωγική συμπεριφορά,
όπου

\[Z' = |Z|\cos t\quad\text{και}\quad Z'' = |Z|\sin t\]

\begin{figure}[H]
  \centering
  \includegraphics[width=0.64\textwidth]{nyquist-full.png}
  \caption{Διάγραμμα Nyquist για δείγμα με gwc 10\% από 20Hz έως 1,5MHz.}
  \label{fig:nyquist-full}
\end{figure}

Η επίδραση των ηλεκτροδίων στα διαγράμματα Nyquist εμφανίζεται ως μια
χαρακτηριστική ουρά 45 μοιρών στις χαμηλές συχνότητες (δεξιά), γεγονός
που επιβεβαιώνεται και στο δείγμα όπως φαίνεται στο \figref{fig:nyquist-full}. Στις
μεσαίες συχνότητες, το φανταστικό μέρος παραμένει σχεδόν σταθερό ενώ το
πραγματικό μειώνεται. Στις υψηλότερες συχνότητες και συγκεκριμένα γύρω
από τη συχνότητα 1MHz και υψηλότερα, παρατηρείται μία απότομη ανοδική
αύξηση στο Z'\,' με σχετική σταθεροποίηση του Ζ'. Αυτό μπορεί να
δικαιολογηθεί από την μετάβαση της ιοντικής αγωγιμότητας που πλέον δε
δύναται να ακολουθήσει το πεδίο, στην επικράτηση της διηλεκτρικής
πόλωσης των μορίων, και κυρίως των διπόλων του νερού, που έχουν υψηλή
διηλετρική σταθερά. Στο \figref{fig:nyquist-trimmed} φαίνεται ευκρινέστερα η απόκριση του
δείγματος μετά την αφαίρεση των συχνοτήτων όπου κυριαρχεί η πόλωση των
ηλεκτροδίων.

Στο \figref{fig:nyquist-moisture} φαίνεται το διάγραμμα Nyquist για τρία ποσοστά υγρασίας
στο ίδιο δείγμα, με gwc 7\%, 10\% και 16\% αντίστοιχα.

\begin{figure}[H]
  \centering
  \includegraphics[width=0.64\textwidth]{nyquist-trimmed.png}
  \caption{Διάγραμμα Nyquist για δείγμα με gwc 10\% εξαιρώντας τα πολωτικά φαινόμενα.}
  \label{fig:nyquist-trimmed}
\end{figure}

\begin{figure}[H]
  \centering
  \includegraphics[width=0.64\textwidth]{nyquist-moisture-comparison.png}
  \caption{Διάγραμμα Nyquist για δείγματα με gwc 7\%, 10\% και 16\%.}
  \label{fig:nyquist-moisture}
\end{figure}

Είναι φανερή η επίδραση της υγρασίας στην απόκριση του συστήματος, αφού
τόσο η πραγματική όσο και η φανταστική συνιστώσα μειώνονται με την
προσθήκη νερού. Για χαμηλά ποσοστά υγρασίας η απόκριση είναι έντονα
εξαρτώμενη από τη συχνότητα, καθώς επηρεάζεται από τη χωρητική
συμπεριφορά του δείγματος. Με την προσθήκη νερού, αυξάνονται τα
διαλυμένα ιόντα και διευκολύνεται η κινητικότητά τους με αποτέλεσμα να
αυξάνεται σημαντικά η αγωγιμότητα. Ως επακόλουθο, η συχνότητα επηρεάζει
σημαντικά λιγότερο την απόκριση.

\section{Ηλεκτρική Μοντελοποίηση του Εδάφους}

Για την καλύτερη κατανόηση των ιδιοτήτων και της συμπεριφοράς ενός μέσου
είναι αναγκαία η μοντελοποίηση του με ένα αντίστοιχο ηλεκτρικό
κυκλωματικό ισοδύναμο. Η συμπεριφορά του εδάφους ωστόσο δεν είναι
απόλυτη, καθώς εξαρτάται από το είδος και το μέγεθος των κόκκων, το
ποσοστό υγρασίας, αλλά και τη συχνότητα στην οποία το εξετάζουμε. Έτσι,
στη βιβλιογραφία αναφέρεται μία πληθώρα κυκλωμάτων ανάλογα με την
εκάστοτε περίπτωση. Έχουν επιλεχθεί να εξεταστούν δύο από τα πιο βασικά
μοντέλα, τα οποία φαίνονται στο \figref{fig:equivalent-model-1} και στο \figref{fig:equivalent-model-2}.

\begin{figure}[H]
  \centering
  \includegraphics[width=0.23\textwidth]{equivalent-model-1.png}
  \caption{Ηλεκτρικό κυκλωματικό ισοδύναμο εδάφους για το μοντέλο 1: παράλληλος συνδυασμός RC.}
  \label{fig:equivalent-model-1}
\end{figure}

\begin{figure}[H]
  \centering
  \includegraphics[width=0.27\textwidth]{equivalent-model-2.png}
  \caption{Ηλεκτρικό κυκλωματικό ισοδύναμο εδάφους για το μοντέλο 2: αντίσταση Rs σε σειρά με παράλληλο συνδυασμό RC.}
  \label{fig:equivalent-model-2}
\end{figure}

Επειδή, όπως αναφέρθηκε, το ισοδύναμο μοντέλο δύναται να αλλάζει ανάλογα
με τη συχνότητα, η σύγκριση έγινε στις εξής δύο περιοχές:

α) Στη χαμηλή ζώνη (100kHz-200kHz) συχνοτήτων όπου μελετάται συνήθως η
αγώγιμη συμπεριφορά, παρακάμπτοντας τα πολωτικά φαινόμενα των
ηλεκτροδίων

β) Στην υψηλή ζώνη (1MHz-1,3MHz) συχνοτήτων όπου μελετάται συνήθως η
χωρητική συμπεριφορά

Για την αξιολόγηση της ακρίβειας των θεωρητικών μοντέλων συγκριτικά με
τα πραγματικά δεδομένα έχει υπολογιστεί το Μέσο Τετραγωνικό Σφάλμα του
κάθε μοντέλου στις δύο παραπάνω περιοχές για μία μέση τιμή υγρασίας
(s10\%) (\tabref{tab:model-mse}).

\begin{table}[H]
  \centering
  \begin{tabularx}{\textwidth}{>{\centering\arraybackslash}X *{4}{>{\centering\arraybackslash}X}}
    \toprule
      & \multicolumn{2}{c}{Z MSE} & \multicolumn{2}{c}{Z Phase} \\
    \cmidrule(lr){2-3}\cmidrule(lr){4-5}
    \textbf{Ζώνη συχνοτήτων} & \textbf{Μοντέλο 1} & \textbf{Μοντέλο 2} & \textbf{Μοντέλο 1} & \textbf{Μοντέλο 2} \\
    \midrule
    Χαμηλή ζώνη συχνοτήτων & 2,27 & 0,72 & 0,0033 & 0,02 \\
    Υψηλή ζώνη συχνοτήτων & 0,048 & 0,18 & 0,00012 & 0,00017 \\
    \bottomrule
  \end{tabularx}
  \caption{Μέσω Τετραγωνικό Σφάλμα για το μέτρο της σύνθετης αντίστασης και της φάσης για δείγμα με gwc 10\% στις δύο περιοχές συχνοτήτων}
  \label{tab:model-mse}
\end{table}

Παρατηρείται πως στις χαμηλότερες συχνότητες το μοντέλο 2 προσομοιώνει
καλύτερα τη συμπεριφορά του δείγματος όσον αφορά το μέτρο της σύνθετης
αντίστασης, ενώ για τη φάση παρά τη μικρή τιμή του μέσου τετραγωνικού
σφάλματος ισχύει το αντίστροφο. Στις υψηλότερες συχνότητες το μοντέλο 1
είναι το βέλτιστο για τη προσομοίωση του μέτρου της σύνθετης αντίστασης,
ενώ η διαφορά των δύο σφαλμάτων για τη φάση είναι εξαιρετικά μικρή.

Επομένως, το δείγμα δεν έχει μία δεδομένη συμπεριφορά, αλλά μπορεί να
μεταβάλλεται στο φάσμα της συχνότητας. Για την αξιοποίηση ενός μοντέλου
είναι απαραίτητη η αξιολόγησή του σε ένα στενό φάσμα κοντά στη συχνότητα
που θα εξεταστεί.

Παρακάτω φαίνεται η έντονη πόλωση των ηλεκτροδίων από τον αναλυτή στην
περιοχή 20Hz-1kHz, καθώς και η απόκριση του δείγματος συγκριτικά με τα
θεωρητικά μοντέλα στις δύο περιοχές συχνοτήτων.

\begin{figure}[H]
  \centering
  \includegraphics[width=0.79\textwidth]{electrode-polarization-analyzer-20hz-1khz.png}
  \caption{Η απόκριση του μέτρου της εμπέδησης και της φάσης σε χαμηλές συχνότητες από 20Hz έως 1kHz με εμφανή την πόλωση των ηλεκτροδίων.}
  \label{fig:electrode-polarization-analyzer}
\end{figure}

\clearpage
\begin{figure}[H]
  \centering
  \includegraphics[width=0.80\textwidth]{model-1-fit-100-200khz.png}
  \caption{Πραγματική και θεωρητική μεταβολή του μέτρου της εμπέδησης και της φάσης για το μοντέλο 1 από 100kHz έως 200kHz.}
  \label{fig:model1-low-fit}
\end{figure}

\begin{figure}[H]
  \centering
  \includegraphics[width=0.80\textwidth]{model-2-fit-100-200khz.png}
  \caption{Πραγματική και θεωρητική μεταβολή του μέτρου της εμπέδησης και της φάσης για το μοντέλο 2 από 100kHz έως 200kHz.}
  \label{fig:model2-low-fit}
\end{figure}

\clearpage
\begin{figure}[H]
  \centering
  \includegraphics[width=0.80\textwidth]{model-1-fit-1-1p5mhz.png}
  \caption{Πραγματική και θεωρητική μεταβολή του μέτρου της εμπέδησης και της φάσης για το μοντέλο 1 από 1MHz έως 1.5MHz.}
  \label{fig:model1-high-fit}
\end{figure}

\begin{figure}[H]
  \centering
  \includegraphics[width=0.80\textwidth]{model-2-fit-1-1p5mhz.png}
  \caption{Πραγματική και θεωρητική μεταβολή του μέτρου της εμπέδησης και της φάσης για το μοντέλο 2 από 1MHz έως 1.5MHz.}
  \label{fig:model2-high-fit}
\end{figure}
